\section{Limitations and Future Work}\label{sec:limits}
\textbf{Lessons.} (1)~\emph{A random split measures memory}; only held-out numbers are evidence \cite{kapoor2023leakage}. (2)~\emph{Train/serve skew hides in plain sight}: raw-$z$ training angles versus served zero-$z$ angles ($14^\circ$ median gap), and a sequence model fed at 0.5\,fps that was trained at 25\,fps, both of which looked like ``model problems'' \cite{sculley2015debt}. (3)~\emph{Silent fallbacks are the dangerous kind}: a retired language-model name returned the safe template for every request and looked like success, and a phone with a faulty GPU path crashed inside a driver instead of ``having no WebGL''. (4)~\emph{No method wins on every photo source}, so report more than one held-out set.

\textbf{Limitations.}
\begin{itemize}
\item \emph{No live-person accuracy yet.} All end-to-end tests use photos or recorded video; the one live test with a person had nobody in view. A study with practitioners and instructors is needed.
\item \emph{The form score is not validated against human judgement.} It is checked only by a perturbation test, it is nearly blind to arm errors (0.94$\to$0.92 when an elbow is broken), and the learned deviation head is close to chance (12.4\% vs.\ 6.7\%); joint-level cues rest on hand-set angle bands.
\item \emph{Weak poses.} Mountain, cobra and child's pose fail on the phone photo test; chair and corpse pass at only 50--60\%; Seated easy failed on one live clip (named upward dog throughout); the ST-GCN reliably handles three poses.
\item \emph{Small data.} 24 videos and public photo sets; some per-pose test sets have fewer than 30 items; single-seed baselines; the held-out public photos share their sources with the training photos (exact near-duplicates are rare, Table~\ref{tab:audit}, but the ranking of methods changes on wild photos, Table~\ref{tab:baseline}); the wild set has only 103 photos.
\item \emph{Sequence-model ablation missing.} The controlled comparison of the ST-GCN with recurrent and convolutional baselines and with a graph-free ST-GCN (Section~\ref{sec:results}) has not been run.
\item \emph{Occlusion.} We abstain rather than recover; the pose \emph{name} can still be influenced by a hallucinated hidden joint; the CLIFF study is small (141 cases, one asymmetric pose, a synthetic box, MediaPipe as reference).
\item \emph{Personalisation.} The range-of-motion profile is not user-tested and a profile learned from a poor stance could forgive real errors; the demonstration profile came from a still photo.
\item \emph{Languages.} Hindi and Bengali strings and paraphrases have not been reviewed by native-speaker instructors, and the forbidden-word screen is a list of \emph{English} words, so it cannot flag unsafe wording in a Hindi or Bengali paraphrase; those languages rely on reviewed templates and the escalation back-off.
\item \emph{Hardware.} Verified on one phone model and desktop browsers; other phones with the same GPU, other browsers, battery and thermal behaviour were not measured. The free-tier server can sleep, in which case basic mode applies (rules only: 36.9\% accuracy and 78.3\% false alarms on held-out photos).
\item \emph{Not a medical device.} We make no claim about injury prevention; no outcome study exists.
\end{itemize}

\textbf{Future work.} (i)~A live user study in which instructors rate form, producing the labelled good/bad-form data the form head lacks; (ii)~retraining the form and deviation heads with arm perturbations and with those labels; (iii)~a tree--network ensemble for naming, since a random forest is the strongest method on same-source photos and the weakest on wild ones, so the two may be complementary (untested); (iv)~viewpoint-robust features (3-D lifting) and a larger pose vocabulary; (v)~opt-in server-side mesh recovery for hidden joints with explicit image consent, now that its benefit for legs and asymmetric poses is measured; (vi)~native-speaker review and language-specific safety screens; (vii)~the sequence-model ablation above and named transitions in the sequence model; (viii)~more devices, including battery and thermal tests.

\section{Conclusion}
AsanaAI shows that a yoga coach can run in a browser, keep the camera image on the device, and still be useful on a budget phone, provided evidence is held to a strict standard. Evaluated only on data never trained on, a gated cascade of residual MLPs names poses on 78.9\% of held-out photos with 20.8\% false alarms and seven passing poses, an ST-GCN with the right input rate reaches 78.0\% macro recall on held-out videos, controlled baselines show that no frame-level method wins on every photo source, and the system degrades gracefully: it abstains on hidden joints, falls back through three pose engines, and keeps coaching offline. We also report what it cannot do yet, which is as important to a person who might follow its advice.

\textbf{Availability.} App: \url{https://yoga-posture-correction-system.vercel.app} (landing page \texttt{/landing}); API: \url{https://arko007-yoga-pose.hf.space}; code: \url{https://github.com/Anamitra-Sarkar/yoga-posture-correction-system}; training data and code: Hugging Face dataset \texttt{Arko007/Yoga-1M} (the transcripts are withheld because they are copyrighted speech).
